\chapter{Optics Crisis: The $f=1$\,mm Design}

\section{Why the Crisis Happened}

The first FD2NN branch inherited an aggressive dual-2f optical layout with
$f=1$\,mm relay lenses. The early numbers looked promising enough to continue,
but the geometry itself was later found to be physically untrustworthy for a
deployable receiver. The issue was not a single failed training run; it was the
sampling and optical scale of the entire design.

\section{Fourier-Plane Sampling Diagnosis}

The key scaling relation is
\begin{equation}
\Delta x_f = \frac{\lambda f}{N \Delta x},
\end{equation}
so the Fourier-plane window width is
\begin{equation}
W_f = N \Delta x_f = \frac{\lambda f}{\Delta x}.
\end{equation}
Shrinking $f$ compresses the Fourier plane and drives the required feature size
into an impractical regime. A complementary aperture argument uses
\begin{equation}
\mathrm{NA} \approx \frac{D}{2f},
\end{equation}
which shows that very small focal lengths raise numerical aperture, clipping
pressure, and fabrication difficulty simultaneously.

\begin{table}[H]
  \centering
  \begin{tabular}{lccc}
    \toprule
    \smalltablehead{Parameter} & \smalltablehead{$f=1$\,mm} & \smalltablehead{$f=10$\,mm} & \smalltablehead{$f=25$\,mm} \\
    \midrule
    $\Delta x_f$ & $\approx 1.5\,\mu$m & $\approx 7.6\,\mu$m & $\approx 18.9$--$37.8\,\mu$m \\
    Approx.\ NA & high ($\sim 0.16$) & moderate ($\sim 0.10$) & lower ($\sim 0.04$) \\
    Diagnostic throughput & $\sim 0.02$ & intermediate & $\sim 0.76$--$1.00$ \\
    Fabrication posture & e-beam scale & difficult & photolithography-friendly \\
    Receiver interpretation & untrustworthy & transitional & physically credible \\
    \bottomrule
  \end{tabular}
  \caption{Conservative scaling comparison for the FD2NN optics crisis. The
  exact NA depends on aperture assumptions, but the trend is robust: $f=1$\,mm
  compresses the Fourier plane too aggressively.}
\end{table}

\section{What the Reinterpretation Means}

The strongest safe conclusion is not that every early FD2NN result became
numerically useless. It is narrower: the $f=1$\,mm family is not a credible
design basis for a deployable optical receiver, because the Fourier-plane
sampling and clipping burden are too severe. This reframing matters because the
later $f=25$\,mm 4f sweep should be read as a corrected design study, not as a
simple continuation of the original branch.

\begin{keyinsight}
The optics crisis corrected the geometry, not the underlying hypothesis. Even
after moving to more credible focal lengths, the campaign still had to answer a
deeper question: can any static phase-only mask improve field-level similarity
for random turbulence?
\end{keyinsight}

That deeper question leads directly to the static-limit analysis in the next
chapter.
